% Options for packages loaded elsewhere
\PassOptionsToPackage{unicode}{hyperref}
\PassOptionsToPackage{hyphens}{url}
\documentclass[
  ignorenonframetext,
]{beamer}
\newif\ifbibliography
\usepackage{pgfpages}
\setbeamertemplate{caption}[numbered]
\setbeamertemplate{caption label separator}{: }
\setbeamercolor{caption name}{fg=normal text.fg}
\beamertemplatenavigationsymbolsempty
% remove section numbering
\setbeamertemplate{part page}{
  \centering
  \begin{beamercolorbox}[sep=16pt,center]{part title}
    \usebeamerfont{part title}\insertpart\par
  \end{beamercolorbox}
}
\setbeamertemplate{section page}{
  \centering
  \begin{beamercolorbox}[sep=12pt,center]{section title}
    \usebeamerfont{section title}\insertsection\par
  \end{beamercolorbox}
}
\setbeamertemplate{subsection page}{
  \centering
  \begin{beamercolorbox}[sep=8pt,center]{subsection title}
    \usebeamerfont{subsection title}\insertsubsection\par
  \end{beamercolorbox}
}
% Prevent slide breaks in the middle of a paragraph
\widowpenalties 1 10000
\raggedbottom
\AtBeginPart{
  \frame{\partpage}
}
\AtBeginSection{
  \ifbibliography
  \else
    \frame{\sectionpage}
  \fi
}
\AtBeginSubsection{
  \frame{\subsectionpage}
}
\usepackage{iftex}
\ifPDFTeX
  \usepackage[T1]{fontenc}
  \usepackage[utf8]{inputenc}
  \usepackage{textcomp} % provide euro and other symbols
\else % if luatex or xetex
  \usepackage{unicode-math} % this also loads fontspec
  \defaultfontfeatures{Scale=MatchLowercase}
  \defaultfontfeatures[\rmfamily]{Ligatures=TeX,Scale=1}
\fi
\usepackage{lmodern}
\ifPDFTeX\else
  % xetex/luatex font selection
\fi
% Use upquote if available, for straight quotes in verbatim environments
\IfFileExists{upquote.sty}{\usepackage{upquote}}{}
\IfFileExists{microtype.sty}{% use microtype if available
  \usepackage[]{microtype}
  \UseMicrotypeSet[protrusion]{basicmath} % disable protrusion for tt fonts
}{}
\makeatletter
\@ifundefined{KOMAClassName}{% if non-KOMA class
  \IfFileExists{parskip.sty}{%
    \usepackage{parskip}
  }{% else
    \setlength{\parindent}{0pt}
    \setlength{\parskip}{6pt plus 2pt minus 1pt}}
}{% if KOMA class
  \KOMAoptions{parskip=half}}
\makeatother
\setlength{\emergencystretch}{3em} % prevent overfull lines
\providecommand{\tightlist}{%
  \setlength{\itemsep}{0pt}\setlength{\parskip}{0pt}}
% Slide layout defaults for warning-clean scientific decks.
\usepackage{etoolbox}
\IfFileExists{xurl.sty}{\usepackage{xurl}}{}
\IfFileExists{seqsplit.sty}{\usepackage{seqsplit}}{\newcommand{\seqsplit}[1]{#1}}
\protected\def\breakseq#1{\seqsplit{#1}}
\protected\def\breaktt#1{\begingroup\ttfamily\seqsplit{#1}\endgroup}
\setlength{\emergencystretch}{6em}
\tolerance=5000
\hbadness=10000
\hfuzz=1pt
\setlength{\tabcolsep}{2pt}
\AtBeginEnvironment{longtable}{\tiny\renewcommand{\arraystretch}{0.86}\setlength{\tabcolsep}{1pt}}
\AtBeginEnvironment{tabular}{\tiny\renewcommand{\arraystretch}{0.86}\setlength{\tabcolsep}{1pt}}

% Natbib and cross-reference fallbacks — slides don't load natbib
% or cleveref, but combined-PDF manuscript prose may emit these
% commands. The fallback renders citations as a bracketed key list
% and cross-references as detokenized labels so slides don't fail on
% undefined control sequences or raw underscores. \providecommand is
% a no-op if packages load later.
\providecommand{\citep}[1]{[#1]}
\providecommand{\citet}[1]{#1}
\providecommand{\citealp}[1]{#1}
\providecommand{\citeauthor}[1]{#1}
\providecommand{\citeyear}[1]{#1}
\providecommand{\cref}[1]{\texttt{\detokenize{#1}}}
\providecommand{\Cref}[1]{\texttt{\detokenize{#1}}}

\newtheorem{proposition}{Proposition}
\newtheorem{hypothesis}{Hypothesis}
\usepackage{bookmark}
\IfFileExists{xurl.sty}{\usepackage{xurl}}{} % add URL line breaks if available
\urlstyle{same}
% fallback for those not using the hyperref driver hyperxmp:
\makeatletter
\@ifundefined{xmpquote}{\newcommand{\xmpquote}[1]{#1}}{}
\makeatother
\hypersetup{
  hidelinks,
  pdfcreator={LaTeX via pandoc}}

\author{\texorpdfstring{}{}}
\date{}

\begin{document}

\section{Discussion}\label{discussion}

\begin{frame}[allowframebreaks]{What the Template Is, and Is Not}
\protect\phantomsection\label{what-the-template-is-and-is-not}
The pipeline measures the \emph{shape} of a literature --- its size,
growth, subfield composition, topical structure, citation geometry, and
the hypotheses a field frames. It does not adjudicate scientific truth.
The optional hypothesis scores summarize how the retrieved corpus
\emph{talks about} each claim, weighted by citation influence; they are
an evidence-landscape instrument, not a verdict.

The 5 topics extracted by NMF provide a data-driven complement to the
keyword-based subfield taxonomy. Where the taxonomy assigns each paper
to a single bucket, the topics reveal overlapping thematic structure: a
paper on modafinil's cognitive effects in ADHD patients belongs to the
``Psychiatry'' subfield but also loads on the ``Cognitive Enhancement''
and ``ADHD Treatment'' topics. This multi-resolution view is more
informative than either approach alone.
\end{frame}

\begin{frame}[allowframebreaks]{Engine Coverage and Bias}
\protect\phantomsection\label{engine-coverage-and-bias}
For this live run, the corpus is dominated by OpenAlex (1,000 records)
and Crossref (1,000 records), with PubMed contributing 986 records and
arXiv contributing 1. Semantic Scholar was rate-limited (HTTP 429) and
returned zero records --- a known limitation of its unauthenticated API
tier. SovietRxiv and ChinaRxiv returned zero records for the modafinil
query, which is expected given their coverage domains.

The max-results cap of 1,000 per engine means the full literature is
larger than the retrieved corpus; the 2302 records represent a bounded
sample rather than the complete literature. The citation network
resolution rate of 22.6\% reflects this: many cited works lie outside
the retrieved slice. Increasing the cap or adding more engines would
improve coverage but also increase runtime and API load.
\end{frame}

\begin{frame}[allowframebreaks]{Honest Defaults}
\protect\phantomsection\label{honest-defaults}
The committed seed corpus is synthetic (reserved test DOIs, generated
authors) so that the whole pipeline runs offline and byte-identically.
Its numbers demonstrate the machinery; they are not empirical findings
about modafinil. Real claims require a live retrieval run with
regenerated figures, reports, and manuscript variables --- as produced
in this instance.
\end{frame}

\begin{frame}[fragile,allowframebreaks]{Limitations and Extensions}
\protect\phantomsection\label{limitations-and-extensions}
Several limitations bound the interpretation of results:

\begin{itemize}
\item
  \textbf{Coverage is bounded by the enabled engines and the query.} The
  max-results cap truncates each engine's contribution. Semantic
  Scholar's rate limiting excluded a major source; a Semantic Scholar
  API key would resolve this.
\item
  \textbf{Subfield classification is keyword-based} and only as good as
  the configured taxonomy. Ambiguous papers may be misclassified; a
  classifier based on embeddings or supervised learning could improve
  accuracy.
\item
  \textbf{The default embeddings are lexical} (TF-IDF/SVD). They capture
  term co-occurrence but not semantic similarity; a transformer backend
  (\texttt{embeddings} extra) would improve the quality of
  nearest-neighbour retrieval and clustering.
\item
  \textbf{Hypothesis scoring depends on an external language model.}
  Without Ollama running, scores read \emph{pending}. The scoring is
  also sensitive to prompt design and model choice; the default
  \texttt{gemma3:4b} is a lightweight model suitable for demonstration
  but may miss nuanced assertions.
\item
  \textbf{Abstract coverage is 55.5\%.} Text analytics operate only on
  the subset of records with abstracts, biasing topic models and
  embeddings toward well-indexed sources.
\end{itemize}

Each limitation is a configuration or dependency choice rather than a
change to the core architecture.
\end{frame}

\end{document}
